\documentclass[10pt,conference,a4paper]{IEEEtran}
\pdfminorversion=6

\usepackage[T1]{fontenc}
\usepackage{cite}
\usepackage{amsmath,amssymb,amsfonts}
\usepackage{graphicx}
\usepackage{textcomp}
\usepackage{booktabs}
\usepackage{multirow}
\usepackage{url}
\usepackage[stretch=50,shrink=50,step=1]{microtype}
\setlength{\columnsep}{0.2in}
\setlength{\textfloatsep}{4pt plus 2pt minus 2pt}
\setlength{\floatsep}{4pt plus 2pt minus 2pt}
\setlength{\intextsep}{4pt plus 2pt minus 2pt}
\setlength{\abovecaptionskip}{3pt}
\setlength{\belowcaptionskip}{2pt}

\newcommand{\sect}[1]{\S#1}

\begin{document}

\title{Post-Hoc Calibration of AI-Content Estimates\\Across Document Length}

\author{
\IEEEauthorblockN{
Pham Quang Dai Phuc,
Pham Thi Minh Hien,
Nguyen Dinh Thuan$^{\ddagger}$
}
\IEEEauthorblockA{
University of Information Technology, Ho Chi Minh City, Vietnam\\
Vietnam National University, Ho Chi Minh City, Vietnam\\
22521137@gm.uit.edu.vn \quad
22520422@gm.uit.edu.vn \quad
thuannd@uit.edu.vn \\
\textsuperscript{$\ddagger$} Corresponding author
}
}


\maketitle

\begin{abstract}
As AI-assisted writing spreads, the useful question is what proportion of a document is
AI-generated --- the percentage Turnitin and Pangram report. Without retraining any
detector, we study how segment scores from 16 existing checkpoints are aggregated into
two outputs: an \emph{estimate} $\hat{\pi}$ of how far a document resembles AI writing
(not the proportion an AI typed), and a \emph{decision} to flag it at a fixed
false-positive rate (FPR). Because segment scores are strongly skewed, both degrade with
document length $k$ when built on the raw score scale. On arXiv Computer Science text,
$\hat{\pi}$ is accurately centered out of sample, but its nominal 95\% Gaussian interval
covers only 86.7\% by $k{=}20$; a distribution-free interval holds coverage if
calibrated near the reported value, but is wide.
Raw-scale thresholds miss a 1\% FPR target in both directions, whereas converting scores
to $z$-scores against human text \emph{before} correcting for within-document
correlation holds FPR near target at every length tested, calibrating on human text
alone; the order matters, as the correlation more than doubles on the transformed
scale, though the gain over the best raw-scale rule is modest. FPR is measured on
full-length human papers; power on full-length AI documents remains unmeasured.
\end{abstract}

\begin{IEEEkeywords}
post-hoc calibration, AI-generated text detection, mixture proportion estimation,
design-effect correction, within-document correlation, document length
\end{IEEEkeywords}

\section{Introduction}

As AI-assisted writing spreads, how much of a document is AI-generated is more useful
than whether it ``contains AI,'' and cannot be answered by simply averaging segment
scores.

We neither train nor modify detectors. We convert their segment scores into a
document-level estimate and decision that remain calibrated with length
(Fig.~\ref{fig:pipe}); the 16 checkpoints of \sect{III-B} are inputs, not contributions.

One claim unifies both results: raw-scale intervals and thresholds fail across document
length because they impose a Gaussian shape on skewed segment scores. Although more
segments should make the mean more Gaussian, neither result improves on real documents
(\sect{V-A}).

\subsection{Background}
Turnitin and Pangram report text percentages; GPTZero instead reports a document-level
authorship probability. These appear only in product documentation, and no examined
public corpus supports evaluation at submission length: median 5{,}560 words
(19 chunks; upper quartile 27; \sect{III-C}). OpAI-Bench (median 244--367 words) and ARB
are about one chunk; RAID~\cite{raid2024} and M4GT-Bench~\cite{m4gt2024} cannot be
grouped into documents. DAGPap24~\cite{dagpap24} is long (5{,}400 words) but leaks
construction artifacts: a 22-feature model using no words reaches AUROC 0.9991.

\subsection{Contributions}
\begin{enumerate}
\item \textbf{An interval for $\hat{\pi}$ across length.} We retain the estimator
(\sect{II-B}) but add an interval for an individual-document estimate.
On real documents (\sect{IV-B}), clustering correction is insufficient; only a
distribution-free interval maintains coverage across $k$. Error skew reverses with
$\pi$, so human calibration covers only $67.2\%$ at $\pi{=}1$ and requires plug-in
recalibration at $\hat{\pi}$.

\item \textbf{Correlation depends on aggregation scale.} With one estimator,
$\rho_{\text{raw}}=0.13$ but $\rho_z=0.30$ (\sect{V-B}). Raw Gaussian and empirical
quantiles miss a 1\% target in opposite directions; transformed aggregation holds
0.76--0.97\%, closer to target though only by $0.44$ points ($p=0.044$). Correlation
remains
essential (18.3\% without it): transformation relocates rather than replaces the design
effect (\sect{V-A}). The ingredients are standard; their ordering is consequential.
\end{enumerate}

\noindent\textbf{Procedure in brief.} The two outputs answer different questions and
never feed each other: the \emph{estimate}, $\hat{\pi}$ with an interval, says how
AI-like a document is; the \emph{decision} says whether to flag it, at an FPR fixed at
any length. Calibration is done once per detector: (i)~human and AI mean segment scores
$\mu_0,\mu_1$ from labeled validation data; (ii)~from about 200 human documents, a pool
of human segment scores and the correlation $\rho_z$ of their transformed scores
(\sect{V-B}); (iii)~error percentiles of $\hat{\pi}$ on mixtures with AI segments at
several known $\pi$ (\sect{IV-B}). For a new document of $k$ segments, report
$\hat{\pi}$ with the interval calibrated at the $\pi$ nearest to it; to decide,
$z$-score each segment against the human pool and flag if the sum exceeds
$2.33\sqrt{k(1+(k-1)\rho_z)}$ (1\% FPR, \eqref{eq:hartung-threshold}).

Three scope conditions apply. The domain is arXiv Computer Science only, deliberately:
$\rho_z$ reflects writing style, and pooling domains would average it away. FPR is
measured at deployment length, but power needs full-length AI papers and remains future
work. Rates assume a clean human/AI dichotomy; on lightly AI-edited human text a nominal
1\% rule operates nearer 5--10\% at segment level (\sect{VI}).

\section{Related Work}

\subsection{Watermarking and distribution-free FPR control}
WaterSeeker~\cite{waterseeker2025} and Zhao et al.~\cite{zhao2025acl} localize
watermarked spans analytically, but dependent overlapping windows preclude deriving
document FPR from window FPR. Multiscale conformal prediction~\cite{mcp2025acl} instead
selects thresholds by length bin; \sect{V-D} adapts it as a baseline. We study empirical
aggregation for learned detectors without analytical nulls, where score skew makes the
apparent intra-document correlation scale-dependent (\sect{V-B}).

\subsection{Estimating AI content proportions at the document or corpus level}
Class-proportion estimation from classifier outputs is \emph{quantification
learning}~\cite{forman2008,hopkinsking2010,gonzalez2017}. PACC~\cite{bella2010} linearly
anchors the two class mean scores; Forman instead corrects hard-positive counts by TPR
and FPR. We retain PACC (\sect{III-D}) and add its document-level interval. Liang et
al.~\cite{liang2024monitoring} work at corpus level, avoiding within-document
correlation; HACo-Det~\cite{haco2025} is the only known peer-reviewed individual-document
study. Neither anticipates that scale dependence, nor the length-dependent
miscalibration it causes.

\subsection{Benchmarks and the length gap}
Sci-SpanDet~\cite{scispandet2026} and DAMASHA~\cite{damasha2026} address span reliability
and authorship transitions, not document proportions. OpAI-Bench (arXiv:2606.06481) and
ARB (arXiv:2607.29539) are unreviewed and cited only as motivation; their 15--23$\times$ length gap from real
papers hides the effect of \sect{V}.

\section{Dataset and Setup}

\subsection{Dataset Construction}
This corpus exists to \emph{evaluate} the post-processing of \sect{IV}--\sect{V}; the
detectors trained on it are inputs. Its design matters only in that shortcuts would
inflate every downstream rate.
From each full-text arXiv paper we extract the abstract and two body paragraphs. Each
human segment is paired with the output of one of 21 models from OpenAI, Anthropic,
Google, Meta, Mistral, DeepSeek and Qwen; hashed segment identifiers fix prompt, mode and
temperature ($0.3$--$1.1$). Identical normalization converts typography to ASCII and
removes introductory phrases, closing a shortcut: print punctuation occurs in 37\% of
AI samples but no LaTeX-derived human sample.

\begin{table}[t]
\caption{Data schema for the corpus.}
\label{tab:fields}
\centering
\scriptsize
\setlength{\tabcolsep}{2.5pt}
\begin{tabular}{@{}lll@{}}
\toprule
Field & Values & Why it is there \\
\midrule
paper identifier & arXiv ID & keys the split, and groups \\
                 &          & every document-level statistic \\
region           & abstract, 2 body paragraphs & fixes what is compared \\
label            & human / AI & --- \\
edit degree      & 0 human, 1 polished, & separates wholly AI text from \\
                 & 3 rewrite, 4 de novo & lightly edited human text \\
generator        & 21 models, 7 providers & --- \\
\bottomrule
\end{tabular}
\end{table}

\begin{table}[t]
\caption{Corpus splits, the share of papers assigned to each by hash, and their use.
The names in the first column are used throughout the paper.}
\label{tab:splits}
\centering
\scriptsize
\setlength{\tabcolsep}{2.5pt}
\begin{tabular}{@{}lll@{}}
\toprule
Split & Share & What it is for \\
\midrule
training              & \multirow{2}{*}{70\%} & fitting the detectors \\
validation            &                       & anchors $\mu_0,\mu_1$, nothing else \\
held-out generator    & 15\%                  & human/AI pairs; Google, DeepSeek only here \\
false positive        & 15\%                  & 3{,}335 human, 955 lightly edited \\
full-length corpus    & separate              & 1{,}200 papers, human only (\sect{III-C}) \\
\bottomrule
\end{tabular}
\end{table}

Table~\ref{tab:fields} lists what each row records and Table~\ref{tab:splits} how the
corpus is divided. A rule fitted and scored within one split uses random
\emph{calibration} and \emph{evaluation halves} of its papers; each repetition is a
\emph{resample}, not a new split. Pairs are
sentence-truncated to matched budgets, discarding the 15\%
with
largest symbol, number and proper-noun discrepancies. Surface-model AUROC then falls
from 0.9343 to 0.8909, and to 0.545 on a length/surface-matched subset of the
held-out-generator split (1{,}206 segments), while the 16 detectors retain
0.927--0.965. Surface cues therefore no longer separate the classes, so what
\sect{IV}--\sect{V} measure is post-processing rather than a corpus artifact.

\subsection{Detector inputs: 16 checkpoints and reference baselines}
The 16 inputs are three seeds each of
RoBERTa-large~\cite{roberta2019}, DeBERTa-v3-large~\cite{deberta2023},
XLM-R-large~\cite{xlmr2020} and ModernBERT-large~\cite{modernbert2024}, plus four
GREATER/RADAR adversarial variants~\cite{greater2025,radar2023}. Held-out-generator AUROC is
0.9746--0.9875 and TPR 0.8435--0.9254 at validation-calibrated $\alpha=1\%$; TF-IDF and
surface baselines reach 0.9491 and 0.8715. Four zero-shot detectors score only
0.40--0.53, which we attribute to matching away the length and token cues curvature
scores use; two are sign-reversed. We neither flip nor retune them, and claim no
comparison. Supplementary Tables S1a--S1b give checkpoint metrics, anchors, $\rho_z$ and
document rates.

\subsection{Full-length corpus}
The full-length corpus contains pre-ChatGPT arXiv Computer Science papers from
2010--2015, excluding all other splits. The same 275-word sentence-boundary segmentation retains
every eligible segment: 27{,}119 from 1{,}200 papers (median 19, mean 22.6, maximum 262).
Its human-only labels permit document FPR but not TPR measurement (\sect{VI}).

\begin{figure}[t]
\centering
\includegraphics[width=\columnwidth]{fig2_pipeline.pdf}
\caption{Pipeline for post-processing. The detector is fixed.}
\label{fig:pipe}
% SOURCE: make_fig2.py
\end{figure}

\subsection{Estimator for $\hat{\pi}$}
Let $\bar{s}$ be the raw mean segment score, unlike the transformed mean $Z_k$ of
\sect{V-B}. The document's AI-like fraction is
\begin{equation}
\hat{\pi}=\frac{\bar{s}-\mu_0}{\mu_1-\mu_0},
\end{equation}
where validation-only $\mu_0,\mu_1$ are wholly human/AI means, independent of both test
sets; evaluation anchoring would make centering tautological. Its expectation equals the
share of AI-written segments only when each segment is wholly human or wholly AI;
otherwise, as for lightly edited text (\sect{VI}), $\hat{\pi}$ measures how closely the
document resembles AI writing, not the share an AI typed. A representative
RoBERTa-large checkpoint has $\mu_0=0.0658$, $\mu_1=0.9643$. Because $\hat{\pi}$ is
unbounded, we clip only for display; all results use it unclipped.

\section{Result 1 --- $\hat{\pi}$ Is Correctly Centered, but Its Interval and Its Threshold Both Depend on $k$}

\subsection{$\hat{\pi}$ matches ground truth: out-of-sample measurement}
% SOURCE: output/pihat_sanity_16models.csv (paper_level_pihat.py)
For each of the 16 checkpoints we compute $\hat{\pi}$ for every sample in
the held-out-generator split using $\mu_0,\mu_1$ estimated exclusively from validation,
giving 637 documents per class; means are computed per checkpoint, then averaged, while
medians and percentiles pool all $10{,}192$ checkpoint--document predictions. Mean
$\hat{\pi}$ is $0.007$ for human and
$0.982$ for AI documents (checkpoint ranges $[-0.009,0.021]$, $[0.966,0.995]$): the
estimator is correctly centered at both extremes though the anchors never saw the
evaluation set. The medians, $-0.038$ and $1.045$, deviate from the means in opposite
directions: the AI distribution is left-skewed with a long left tail ($p_5=0.531$), the
human one right-skewed ($p_{95}=0.325$). This is the first evidence for the central
claim of \sect{V}: segment scores are strongly non-Gaussian, and $5\%$ of wholly human
documents already report above $\hat{\pi}=0.32$.

Centering at the endpoints does not establish accuracy in between, so we built
$k{=}20$ documents replacing a controlled random fraction of segments by their AI
counterparts, which that split supplies for the same 1{,}514 slots.
Across $\pi\in\{0,0.1,0.25,0.5,0.75,0.9,1\}$ mean $\hat{\pi}$ is $0.007$, $0.104$,
$0.246$, $0.492$, $0.737$, $0.878$ and $0.977$ (median absolute error $0.039$).
Linearity is a property of an affine estimator, not a finding; what the endpoints
cannot predict is the error's size, almost entirely dispersion rather than bias
(median $|\text{bias}|/\mathrm{SD}=0.15$).

\subsection{An interval for $\hat{\pi}$, and what it costs}
A point estimate needs uncertainty. As a cluster sample, a document's segments give
$\hat{\pi}$ the design effect \sect{V} applies to a threshold:
\begin{equation}
\widehat{\mathrm{SE}}(\hat{\pi})=\frac{1}{\mu_1-\mu_0}
\sqrt{\frac{\hat{\sigma}^2}{k}\bigl(1+(k-1)\rho_{\text{raw}}\bigr)},
\label{eq:pi-se}
\end{equation}
with $\hat{\sigma}^2$ the pooled segment-score variance and $\rho_{\text{raw}}$ the
raw-scale intra-class correlation (median $0.132$, range 0.087--0.167), both estimated
on a calibration half of the full-length corpus, disjoint from the documents measured.
We compare $\hat{\pi}\pm z_{0.975}\widehat{\mathrm{SE}}$, with and without the clustering
factor, against a distribution-free interval taking empirical $2.5$/$97.5$ percentiles
on the same half.

\begin{table}[t]
\caption{Coverage and width of a nominal $95\%$ interval for $\hat{\pi}$ on human
documents.}
\label{tab:pici}
% SOURCE: final_experiments_rivf2026/results/f1_coverage_null.csv (f1_pihat_ci.py)
\centering
\scriptsize
\setlength{\tabcolsep}{3pt}
\begin{tabular}{@{}rccccc@{}}
\toprule
& \multicolumn{3}{c}{Coverage (\%)} & \multicolumn{2}{c}{Half-width} \\
\cmidrule(lr){2-4}\cmidrule(lr){5-6}
$k$ & Gaussian & $+$design eff. & Empirical & Gauss. & Emp. \\
\midrule
1  & 95.7 & 95.7 & 94.8 & 0.368 & 0.493 \\
5  & 90.5 & 93.3 & 94.8 & 0.165 & 0.198 \\
10 & 89.5 & 93.5 & 94.8 & 0.116 & 0.163 \\
20 & 86.7 & 93.8 & 94.8 & 0.082 & 0.152 \\
\bottomrule
\end{tabular}
\end{table}

Table~\ref{tab:pici} gives that coverage, measured on the 576 papers of the full-length corpus with
at least 20 segments --- the controlled population of \sect{V-A}, where within-document
correlation is genuine rather than simulated. The Gaussian interval is correct at
$k{=}1$ and decays to $86.7\%$ at $k{=}20$; the design effect recovers most of that gap
without closing it; only the distribution-free interval holds $94.8\%$ at every $k$,
and it also varies least across checkpoints ($93.5$--$95.6\%$ at $k{=}20$).

The honest interval is the wide one: a full paper admits about $\pm0.15$, so
$\hat{\pi}=0.30$ and $0.55$ are not reliably distinguishable, and more text barely helps
(half-width $0.198$, $0.163$, $0.152$ at $k{=}5,10,20$). That is four times the MAE of
\sect{IV-A}, whose mixtures draw segments from different papers and so lack
within-document correlation: synthetic error understates real uncertainty.

% SOURCE: final_experiments_rivf2026/results/f6_coverage_mixed_emp.csv (f6_mixed_pi_empirical_ci.py)


% >>> TABLE_MIXED_BEGIN -- sinh boi make_table_mixed.py, DUNG sua tay
\begin{table}[t]
\caption{Interval coverage (\%) on documents containing AI: $k{=}20$ human
segments from one real paper, $\pi k$ of them replaced by AI scores.}
\label{tab:mixed}
% SOURCE: make_table_mixed.py <- f6_coverage_mixed_emp.csv (f6_mixed_pi_empirical_ci.py)
\centering
\scriptsize
\setlength{\tabcolsep}{3.2pt}
\begin{tabular}{@{}rccccc@{}}
\toprule
& \multicolumn{3}{c}{Empirical, calibrated at} & \multicolumn{2}{c}{Gaussian} \\
\cmidrule(lr){2-4}\cmidrule(lr){5-6}
$\pi$ & $\pi{=}0$ & $\hat{\pi}$ (plug-in) & true $\pi$ & $+$deff & plain \\
\midrule
0.00 & 94.6 & 93.0 & 94.6 & 94.0 & 87.7 \\
0.25 & 79.6 & 93.2 & 95.1 & 95.9 & 85.7 \\
0.50 & 72.5 & 94.1 & 94.9 & 97.2 & 84.7 \\
0.75 & 76.4 & 93.7 & 94.9 & 98.7 & 86.9 \\
1.00 & 67.2 & 93.8 & 94.7 & 97.4 & 86.4 \\
\midrule
\multicolumn{6}{@{}l@{}}{\footnotesize Plug-in half-width narrows $0.147\rightarrow0.101$ as $\hat{\pi}$ goes $0\rightarrow1$.}\\
\bottomrule
\end{tabular}
\end{table}
% <<< TABLE_MIXED_END

Table~\ref{tab:pici} covers human documents only. Table~\ref{tab:mixed} repeats the
measurement on documents that do contain AI, and gives the cost:
an interval calibrated there collapses once a document contains AI. Width is not the
cause, that interval being the widest at $\pm0.147$; position is. Error skew is
right-tailed under a human null and left-tailed under an AI one, so calibration at one
end misfits the other. As $\pi$ is unknown in deployment, the plug-in $\hat{\pi}$
interval holds throughout. A reported percentage should carry an interval calibrated at
that percentage.

\subsection{Incorrectly set thresholds fail in opposite directions}
% SOURCE: output/mean_ktransfer_16models.csv; ranges -> s4_ranges_main_tables.csv


A decision rule can be miscalibrated even when $\hat{\pi}$ is centered, because the
right threshold depends on the length of the population used to set it.
% >>> TABLE_KTRANSFER_BEGIN -- sinh boi make_table_ktransfer.py, DUNG sua tay
\begin{table}[t]
\caption{Cost of pinning the threshold at the wrong length. Only $k{\approx}2.38$
matches the split it is scored on. Percentages.}
\label{tab:ktransfer}
% SOURCE: make_table_ktransfer.py <- output/mean_ktransfer_16models.csv
\centering
\scriptsize
\setlength{\tabcolsep}{2.5pt}
\begin{tabular}{@{}lrrrrrr@{}}
\toprule
& \multicolumn{2}{c}{$k{\approx}1$} & \multicolumn{2}{c}{$k{\approx}22.6$} & \multicolumn{2}{c}{$k{\approx}2.38$} \\
\cmidrule(lr){2-3}\cmidrule(lr){4-5}\cmidrule(lr){6-7}
Model & FPR & Miss & FPR & Miss & FPR & Miss \\
\midrule
RoBERTa-large (s44) & 0.00 & 16.33 & 3.77 & 2.04 & 1.26 & 6.16 \\
RoBERTa-large (s43) & 0.00 & 16.33 & 4.40 & 1.10 & 1.30 & 3.87 \\
GREATER + RoBERTa-large (s42) & 0.16 & 16.80 & 4.08 & 2.04 & 1.30 & 5.30 \\
XLM-R-large (s42) & 0.16 & 19.15 & 4.87 & 1.10 & 1.32 & 6.22 \\
XLM-R-large (s44) & 0.16 & 19.31 & 5.18 & 1.26 & 1.23 & 10.17 \\
GREATER + XLM-R-large (s42) & 0.16 & 19.78 & 6.59 & 0.94 & 1.24 & 5.56 \\
XLM-R-large (s43) & 0.16 & 19.94 & 4.08 & 0.94 & 1.27 & 6.68 \\
RADAR (s42) & 0.16 & 21.04 & 6.28 & 1.57 & 1.30 & 4.51 \\
DeBERTa-v3-large (s42) & 0.31 & 13.50 & 6.75 & 0.63 & 1.32 & 8.56 \\
DeBERTa-v3-large (s44) & 0.31 & 15.07 & 5.34 & 1.26 & 1.37 & 10.74 \\
ModernBERT-large (s43) & 0.31 & 16.17 & 4.71 & 0.94 & 1.40 & 6.16 \\
ModernBERT-large (s44) & 0.31 & 18.84 & 7.85 & 0.78 & 1.32 & 3.85 \\
DeBERTa-v3-large (s43) & 0.47 & 15.86 & 5.97 & 0.63 & 1.31 & 5.17 \\
RoBERTa-large (s42) & 0.47 & 17.74 & 5.02 & 1.88 & 1.34 & 11.61 \\
ModernBERT-large (s42) & 0.47 & 17.90 & 4.40 & 0.78 & 1.29 & 4.84 \\
RADAR + XLM-R-large (s42) & 0.63 & 24.80 & 5.34 & 4.40 & 1.37 & 20.50 \\
\midrule
\textbf{Median} & \textbf{0.24} & \textbf{17.82} & \textbf{5.10} & \textbf{1.10} & \textbf{1.31} & \textbf{6.16} \\
\bottomrule
\end{tabular}
\end{table}
% <<< TABLE_KTRANSFER_END

Table~\ref{tab:ktransfer} scores three sources for $\tau$ on the same split
(held-out-generator, $k\approx2.38$, 637 pairs): a segment-level threshold from
validation ($k\approx1$), one pinned at $k\approx22.6$ on the full-length corpus, and
one calibrated at the evaluation length itself (fitted on a calibration half, scored on
the other over 200 resamples; the first two on all 637 pairs). All three lie on one
trade-off curve, so the failure is not a worse curve but that pinning $\tau$ at the
wrong $k$ lands far from the $\alpha$ requested, too strict from below and too
permissive from above. That holds for every checkpoint, and the spread is wide enough
that the median understates the risk: the correct-$k$ control misses
$3.85$--$20.50\%$.


\section{Result 2 --- $z$-Transformation First, Correlation Adjustment Later}

\subsection{Why direct Gaussian design-effect adjustment fails on the raw scale}
We use a controlled design over the same 576 documents throughout, those with at least
20 segments --- $48\%$ of the corpus, long enough for the full grid. Only $k$ changes,
by within-document subsampling; no document enters or leaves. Results aggregate across
16 checkpoints and five resamples.

\begin{figure}[t]
\centering
\includegraphics[width=\columnwidth]{fig1_fpr_curve.png}
\caption{Document-level FPR against a $1\%$ target.}
\label{fig:fpr}
% SOURCE: make_fig1.py <- final_experiments_rivf2026/results/f7_cluster_marginal_ci.csv
%         (static, emp_deff, zscale) + output/pihat_static_vs_corrected_by_k.csv (Gauss)
\end{figure}

Both raw-scale rules threshold the document mean $\bar{s}$ at
\begin{equation}
\tau(k)=\mu_0+(q-\mu_0)\sqrt{\frac{1+(k-1)\rho_{\text{raw}}}{k}},
\label{eq:raw-rule}
\end{equation}
differing only in $q$: the Gaussian variant sets $q=\mu_0+z_{1-\alpha}\hat{\sigma}$,
the empirical one the observed $99$th percentile of a single calibration segment. As Fig.~\ref{fig:fpr} shows, the Gaussian variant never brings FPR
near $1\%$: it jumps to $6.78\%$ at $k{=}2$ then settles around $5.4\%$, five times the
target, with no sign of improving as $k$ grows. The factor $1+(k-1)\rho_{\text{raw}}$
corrects only the variance; the threshold built from it still assumes a Gaussian null.
Under that null $\mu_0\approx0.07$ while the right tail reaches nearly $1$, so
$z_{1-\alpha}\sigma$ falls well short of the true $99$th percentile.

That averaging never rescues it follows from \eqref{eq:raw-rule}, whose factor is not
$1/\sqrt{k}$ but $\sqrt{(1+(k-1)\rho)/k}\to\sqrt{\rho}$: $\tau(k)$ flattens to a floor
$\mu_0+(q-\mu_0)\sqrt{\rho}$ set by $\rho$ alone, which is why the curve plateaus. With
$\rho_{\text{raw}}=0.13$ that floor retains $36\%$ of the distance from $\mu_0$ to $q$,
applied to a $q$ already short; on the $z$-scale $\rho_z=0.303$ puts it at $55\%$. The
central limit theorem improves the shape of $\bar{s}$, but not fast enough to overcome a
threshold converging to the wrong place.

% SOURCE: final_experiments_rivf2026/results/f5_deff_necessity.csv (f5_deff_necessity.py)
To isolate the ingredients, replacing the Gaussian quantile by the empirical $99$th
percentile of a single segment, with design effect and raw scale untouched, fixes the
skew but overcorrects: FPR falls to $0.33$--$0.47\%$ for every $k>1$, a third of target.
Both variants scale $q$ by the same Gaussian-motivated shrinkage term of
\eqref{eq:raw-rule}, so choosing $q$ empirically cannot repair it; the two raw-scale
quantiles miss in opposite directions. Deleting the clustering factor instead ($\rho=0$,
the rule otherwise identical) is worse still: on the $z$-scale FPR rises
$0.97\rightarrow5.40\rightarrow18.28\%$ at $k=1,5,20$, on the raw scale
$1.10\rightarrow2.54\%$. The factor is necessary on both scales, and does most work on
the transformed one, where $\rho$ is over twice as large.

At $k{=}1$ the empirical and $z$-scale rules use the same order statistic, so their
residual gap ($1.10$ against $0.97\%$) is quantile-estimation error; every divergence at
$k>1$ traces to the averaging scale, and hence to which $\rho$ governs the design effect.

Since both rules score the same papers, we bootstrap their \emph{difference}
($B{=}2000$), resampling paper clusters because the five evaluation halves overlap. The
$z$-scale rule is closer to target, but only by $0.44$ points of root-mean-square and
$0.43$ of worst-case deviation ($[-0.53,-0.02]$, $p=0.044$; $[-0.58,-0.02]$,
$p=0.047$); both beat the static threshold far more decisively ($p<0.001$). The
empirical rule errs on the safe side, significantly \emph{below} $1\%$ at every $k>1$,
while the $z$-scale rule contains the target at every $k$. This ranks calibration, not
safety; \sect{V-C} adds power at matched FPR.

Every working rule sits slightly \emph{below} target: $T_k=Z_k/\sqrt{1+(k-1)\rho_z}$ is
near-Gaussian but under-dispersed at $k{=}20$ (sd $0.965$), so Hartung's estimator
overstates $\rho_z$ by about $9\%$ and the rule is conservative by construction.

\subsection{Mechanism: transformation before correlation adjustment}
For each segment score $s_i$, we compute a continuity-corrected empirical right-tailed
$p$-value against a pool of human calibration scores $s^{(0)}_1,\ldots,s^{(0)}_m$
held out from the documents evaluated, so a small $p_i$ marks a segment scoring higher
than most human calibration text:
\begin{equation}
\begin{aligned}
p_i&=\frac{\#\{j:s^{(0)}_j\ge s_i\}+\tfrac{1}{2}}{m+1},\\
z_i&=\Phi^{-1}(1-p_i).
\end{aligned}
\label{eq:pit-transform}
\end{equation}
Under $H_0$ (all segments human and exchangeable with the calibration population),
$p_i$ is approximately uniform and $z_i$ therefore approximately $N(0,1)$ by the
probability integral transform, whatever the skew of the original scores. The pool is
always disjoint from the documents it is applied to, holding $m\approx9{,}000$ segments
in \sect{V-A} but only $m{=}438$ in the \sect{V-C} transfer, where just four discrete
$p$-values fall below $\alpha{=}1\%$ --- a transfer limited by tail resolution as much
as by sample size.

On this scale, define the aggregated statistic $Z_k=k^{-1/2}\sum_{i=1}^{k}z_i$.
Write $\rho_z$ for the common within-document correlation \emph{on this transformed
scale}, against the raw-scale $\rho_{\text{raw}}$ of \eqref{eq:pi-se}; under
exchangeability the variance of $Z_k$ is multiplied by $1+(k-1)\rho_z$. The resulting decision rule is
\begin{equation}
Z_k>z_{1-\alpha}\sqrt{1+(k-1)\rho_z},
\label{eq:hartung-threshold}
\end{equation}
where $z_{1-\alpha}=2.33$ for $\alpha=1\%$. We estimate $\rho_z$ per checkpoint by the
Hartung formula~\cite{hartung1999}. Averaging the five resample estimates within each
checkpoint gives median 0.303 and range 0.204--0.401 (the 80 resample-level estimates span
0.183--0.424): $\rho_z$ is a checkpoint-specific calibration quantity, not a
universal constant.

% SOURCE: output/rho_compare_16models.csv (paper_level_rho_compare.py); icc.py
This 0.303 against the $\rho_{\text{raw}}=0.132$ of \eqref{eq:pi-se} is not an artifact
of two estimators: the same one-way ICC on the same documents and $20$ segments each
returns 0.296 on transformed scores, so the estimator moves $\rho$ by under a hundredth
while the scale moves it from 0.132 to 0.296. The correlation a variance-inflation
formula needs belongs to the scale one aggregates on.

$\rho_z$ must be estimated with care. Perturbing it by $\pm20\%$ moves FPR by under a
point but the miss rate on half-AI documents from $17.3\%$ to $48.5\%$ at $k{=}20$. Its
checkpoint range is wider still, so a $\rho_z$ carried to a new checkpoint holds FPR
without predicting power; those miss rates are optimistic anyway (\sect{VI}).

That fixes a calibration budget. Resampling the calibration half, the $95\%$ interval
for $\rho_z$ is $\pm54\%$ of its value at 25 human documents, $\pm27\%$ at 100 and
$\pm17\%$ at 200. Since $\pm20\%$ is the perturbation above, \textbf{roughly 200 human
documents} is the practical floor for adopting this rule on a new detector, and they
need no AI labels: $\rho_z$ is estimated under the null.

The reference point is an oracle recalibrated per $k$ on labeled data at the
deployment length. The paired bootstrap resolves no difference between them: $-0.12$
points of RMS deviation $[-0.20,+0.14]$, $-0.08$ of worst-case $[-0.26,+0.26]$; since
non-rejection is not equivalence, we report the interval. The oracle needs labeled data
at that length; this rule needs none, and its per-checkpoint FPR stays within
$0.31$--$1.24\%$.

\subsection{Out-of-sample generalization}
We estimate $\rho_z$ on the full-length corpus and apply it unchanged to the
held-out-generator split: distinct documents of the same corpus at a natural mean length
$k\approx2.38$, not used to estimate it.

% SOURCE: output/hartung_ood_16models.csv; ranges -> s4_ranges_main_tables.csv


% >>> TABLE_OOD_BEGIN -- sinh boi make_table_ood.py, DUNG sua tay
\begin{table}[t]
\caption{$\rho_z$ from the full-length corpus, reused unchanged on the
held-out-generator split ($k{\approx}2.38$); median [range] over 16 checkpoints.
Rows 2--4 threshold the raw mean: fitted at this length on calibration halves
(Table~\ref{tab:ktransfer}), then scored on $Z_k$'s 637 pairs, then re-fitted to
$Z_k$'s FPR. Wins: checkpoints where $Z_k$ misses fewer.}
\label{tab:ood}
% SOURCE: make_table_ood.py <- hartung_ood_16models.csv, mean_ktransfer_16models.csv,
%          rev3_paired_miss.csv, rev4_fpr_matched.csv
\centering
\scriptsize
\setlength{\tabcolsep}{3pt}
\begin{tabular}{@{}lccc@{}}
\toprule
Rule & FPR (\%) & Miss (\%) & Wins \\
\midrule
1. $Z_k$ (proposed) & 1.26\,[0.47,1.88] & 3.77\,[2.67,8.16] & --- \\
\midrule
2. Raw mean, fitted here & 1.31\,[1.23,1.40] & 6.16\,[3.85,20.50] & --- \\
3. Same, on $Z_k$'s pairs & 1.10 & 6.04 & 13/16 \\
4. Raw mean, FPR of $Z_k$ & \emph{matched} & 5.02 & 14/16 \\
\bottomrule
\end{tabular}
\\[2pt]\footnotesize In row 4, 419 checkpoint--AI-document pairs are flagged only
by $Z_k$ and 90 only by the raw mean; unadjusted McNemar tests ($5\%$) favour
$Z_k$ for 7 checkpoints and the raw mean for 1.
\end{table}
% <<< TABLE_OOD_END

Table~\ref{tab:ood} reports what survives the move. Calibration carries over, at a lower
miss rate than a raw-mean threshold fitted at this very length (rows 1--2). Part of
that gain is a looser operating point, since a lower threshold trades miss against FPR
mechanically (row 3); re-fitting the raw-mean threshold to flag exactly as many human
papers as $Z_k$ roughly halves the advantage but does not remove it (row 4). What
remains is consistent with a change in ranking: the mean of monotonically transformed
scores is not monotone in the raw mean. $\rho_z$ thus carries its \emph{calibration}
to another split; power is a different matter, since \sect{V-B} shows $\rho_z$ does not
predict it.

\subsection{Comparison with per-bin conformal calibration}
% SOURCE: final_experiments_rivf2026/results/f3_conformal_overall.csv (f3_conformal_bins.py)
The competitor from prior work is length-binned split conformal calibration, inspired
by MCP~\cite{mcp2025acl} and run on our own scores at natural length (1{,}200 papers,
median $k=19$, range 10--262). We bin by equal counts of $k$ rather than MCP's
equal-width token intervals, so this adapts rather than reproduces that algorithm. With
five bins it attains median FPR $0.97\%$ $[0.70,1.10]$ against $0.77\%$ $[0.27,1.77]$
for the proposed rule and $1.25\%$ for a static threshold: a strong baseline we do not
claim to beat on FPR. Its cost is resolution. A conformal threshold at level $\alpha$
needs $\lceil 1/\alpha\rceil-1=99$ calibration documents \emph{per bin}, so 600 papers
cap the design at six; at ten every bin falls short and the rule never flags. The
proposed rule uses each document's own $k$, tied to neither $\alpha$ nor the budget.
That, not a lower FPR, is the difference.

\section{Discussion and Limitations}
Human text lightly edited by AI draws segment-level FPRs 4.7--9.8$\times$ higher than
purely human text ($n{=}955$ against $3{,}335$). The two-component model places every
segment at one extreme, whereas polished text lies between --- the reason, noted in
\sect{III-D}, that $\hat{\pi}$ measures resemblance to AI writing, not authorship. As a validity check, a
classifier reading only segment counts achieves AUROC 0.5000.

TPR on \emph{real} AI documents at deployment length remains unmeasured: the full-length
corpus holds only human text, so the miss rates in \sect{IV} and \sect{V-C} come from
the held-out-generator split at $k\approx2$--3, and the $k{=}20$ figures in \sect{V-B}
use documents assembled from different papers, whose AI half lacks within-document
correlation and is thus optimistic. This bounds where the
contribution applies: at $k\ge8$ our synthetic \emph{fully} AI documents are detected
almost without error by every rule tested, so a decision rule earns its keep on short
documents and on mixed $\pi$, not long fully-AI ones. Closing the gap needs full-length
AI papers rather than reassembled segments, and the count is modest: at an expected
miss near $4\%$, about 350 such documents pin it to $\pm2$ points and 1{,}400 to
$\pm1$.

Within arXiv Computer Science, \sect{V-C} shows $\rho_z$ transferring between splits;
we claim no more, and since $\rho_z$ tracks writing style a new domain would need its
own estimate --- a recalibration, not a retraining. Every
result also assumes one segmentation, 275 words at sentence boundaries; since $\rho_z$
is a correlation between segments, a coarser or finer segmentation would move it, and we have
not measured by how much. Relatedly, $k$ in
\sect{V-A} is a subsample size from papers of at least 20 segments, so the $k{=}1$ row
is one segment of a long paper, not a short submission. FPR also varies by subfield,
from $0.00\%$ (\texttt{cs.LG}) to $7.21\%$ (\texttt{cs.PL}) over the 19 with at least 20
papers.

Finally, the human side is arXiv Computer Science 2010--2015, chosen so undisclosed
LLM assistance is implausible, while the AI side comes from 2024--26 models. Topic is matched but period is not, and \sect{III-A} controls punctuation, symbol and
length cues rather than era. The 0.97--0.99 AUROC may thus partly reflect an era
boundary, making miss rates and $\hat{\pi}$ on AI-containing documents look better than they are, though not
the FPR results, which are calibrated and measured on human text of one period.
Deployment needs the human side ($\mu_0$, pool, $\rho_z$) re-estimated on contemporary
writing, without AI labels; otherwise $\mu_0$ may sit below the 2020s human anchor,
biasing $\hat{\pi}$ upward.

\section{Conclusion}
$\hat{\pi}$ is accurately centered out of sample, but neither its interval nor a
decision threshold built on the raw score scale survives a change in document length,
for one reason: a Gaussian shape assumption over skewed segment scores, whose error
compounds with $k$. The fault is that assumption, not the clustering correction, which
stays indispensable; the transform moves the correction onto a scale where $\rho$ is
over twice as large and a Gaussian quantile is already right.

Under a cluster bootstrap the transform is closer to target than the raw-scale
empirical rule, though only marginally ($p=0.044$, a gap of $0.44$ points), and is
indistinguishable from a per-$k$ oracle it never needs, while per-bin
conformal calibration is better centered where the budget allows. What further
distinguishes it is power at matched FPR (\sect{V-C}) and length resolution independent
of that budget (\sect{V-D}). The claim the evidence supports without qualification is
negative: Gaussian intervals and thresholds read off untransformed scores should not be
used when documents vary in length, as papers do. Being post-processing, the remedy
costs no retraining and carries to any scorer exposing per-segment scores, including the
closed systems of \sect{I} should they do so: calibration reads human text only, so FPR
control does not depend on how well the scorer separates. Power does, and we tested only
scorers above AUROC $0.92$, excluding the zero-shot detectors of \sect{III-B}. A
percentage should be reported
with an interval, wide at about $\pm0.15$ and calibrated at that percentage.

\noindent\textbf{Reproducibility.} Code, per-checkpoint breakdowns and every analysis
above are available at \url{https://github.com/hien01604/posthoc-ai-content-calibration}.

\section*{Acknowledgment}

This research was supported by the VNUHCM-University of Information Technology's Scientific Research Support Fund.

\begin{thebibliography}{00}

\bibitem{raid2024}
L. Dugan \textit{et al.}, ``RAID: A Shared Benchmark for Robust Evaluation of Machine-Generated
Text Detectors,'' in \textit{Proc. ACL}, 2024, pp. 12463--12492.

\bibitem{m4gt2024}
Y. Wang \textit{et al.}, ``M4GT-Bench: Evaluation Benchmark for Black-Box Machine-Generated
Text Detection,'' in \textit{Proc. ACL}, 2024, pp. 3964--3992.

\bibitem{dagpap24}
S. Chamezopoulos \textit{et al.}, ``Overview of the DagPap24 Shared Task on Detecting
Automatically Generated Scientific Paper,'' in \textit{Proc. SDP}, 2024, pp. 7--11.

\bibitem{waterseeker2025}
L. Pan \textit{et al.}, ``WaterSeeker: Pioneering Efficient Detection of Watermarked
Segments in Large Documents,'' in \textit{Findings of NAACL}, 2025, pp. 2866--2882.

\bibitem{zhao2025acl}
X. Zhao \textit{et al.}, ``Efficiently Identifying Watermarked Segments in Mixed-Source
Texts,'' in \textit{Proc. ACL}, 2025, pp. 6304--6316.

\bibitem{mcp2025acl}
X. Zhu \textit{et al.}, ``Reliably Bounding False Positives: A Zero-Shot
Machine-Generated Text Detection Framework via Multiscaled Conformal Prediction,'' in
\textit{Proc. ACL}, 2025, pp. 12298--12319.

\bibitem{forman2008}
G. Forman, ``Quantifying Counts and Costs via Classification,'' \textit{Data Mining and
Knowledge Discovery}, vol. 17, no. 2, pp. 164--206, 2008.

\bibitem{hopkinsking2010}
D. J. Hopkins and G. King, ``A Method of Automated Nonparametric Content Analysis for
Social Science,'' \textit{American Journal of Political Science}, vol. 54, no. 1,
pp. 229--247, 2010.

\bibitem{gonzalez2017}
P. Gonz\'alez \textit{et al.}, ``A Review on Quantification Learning,'' \textit{ACM
Computing Surveys}, vol. 50, no. 5, Art. no. 74, 2017.

\bibitem{bella2010}
A. Bella \textit{et al.}, ``Quantification via Probability Estimators,'' in
\textit{Proc. IEEE ICDM}, 2010, pp. 737--742.

\bibitem{liang2024monitoring}
W. Liang \textit{et al.}, ``Monitoring AI-Modified Content at Scale: A Case Study on
the Impact of ChatGPT on AI Conference Peer Reviews,'' in \textit{Proc. ICML}, 2024, pp. 29575--29620.

\bibitem{haco2025}
Z. Su \textit{et al.}, ``HACo-Det: A Study Towards Fine-Grained Machine-Generated Text
Detection under Human--AI Coauthoring,'' in \textit{Proc. ACL}, 2025, pp. 22015--22036.

\bibitem{scispandet2026}
Z. Yin and S. Wang, ``Span-Level Detection of AI-Generated Scientific Text via
Contrastive Learning and Structural Calibration,'' \textit{Knowledge-Based Systems},
vol. 334, Art. no. 115123, 2026.

\bibitem{damasha2026}
L. D. M. S. Sai Teja \textit{et al.}, ``DAMASHA: Detecting AI in Mixed Adversarial Texts
via Segmentation with Human-Interpretable Attribution,'' in \textit{Findings of EACL},
2026, pp. 6189--6206.

\bibitem{roberta2019}
Y. Liu \textit{et al.}, ``RoBERTa: A Robustly Optimized BERT Pretraining
Approach,'' arXiv:1907.11692, 2019.

\bibitem{deberta2023}
P. He \textit{et al.}, ``DeBERTaV3: Improving DeBERTa Using ELECTRA-Style Pre-Training
With Gradient-Disentangled Embedding Sharing,'' in \textit{Proc. ICLR}, 2023.

\bibitem{xlmr2020}
A. Conneau \textit{et al.}, ``Unsupervised Cross-lingual Representation Learning at
Scale,'' in \textit{Proc. ACL}, 2020, pp. 8440--8451.

\bibitem{modernbert2024}
B. Warner \textit{et al.}, ``Smarter, Better, Faster, Longer: A Modern Bidirectional
Encoder for Fast, Memory Efficient, and Long Context Finetuning and Inference,'' in
\textit{Proc. ACL}, 2025, pp. 2526--2547.

\bibitem{greater2025}
Y. Li \textit{et al.}, ``Iron Sharpens Iron: Defending Against Attacks in
Machine-Generated Text Detection With Adversarial Training,'' in \textit{Proc. ACL},
2025, pp. 3091--3113.

\bibitem{radar2023}
X. Hu \textit{et al.}, ``RADAR: Robust AI-Text Detection via Adversarial Learning,'' in
\textit{Proc. NeurIPS}, 2023.

\bibitem{hartung1999}
J. Hartung, ``A Note on Combining Dependent Tests of Significance,'' \textit{Biometrical
Journal}, vol. 41, no. 7, pp. 849--855, 1999.

\end{thebibliography}




\end{document}